%%%%%%%%%%%%%%%%%%%%%%%%%%%%%%%%%
%
%       EXERCÍCIO
%
%%%%%%%%%%%%%%%%%%%%%%%%%%%%%%%%%

\definevalorquestao

\begin{exercicioBanco}[\valorquestao]

Pesquisadores desenvolveram 2 medicamentos diferentes para o tratamento de uma doença. A probabilidade de cura dessa doença caso o paciente não receba nenhum tratamento é 20\%. O medicamento A garante probabilidade de cura igual a 80\% considerando que o paciente não tenha nenhuma outra comorbidade. A cada doença associada, a probabilidade de cura cai 5 pontos percentuais (5\%). O medicamento B garante probabilidade de cura igual a 60\% independentemente da comorbidade do paciente. A partir de quantas comorbidades se torna mais indicado o uso do medicamento B?

\end{exercicioBanco}


